\section{Exact arguments for Fisher's estimation core}
\label{sec:annex-p15-fisher-estimation}

\begin{annexitem}{annex:p15.fe.binomial}{thm:p15-fe-binomial}{The exact binomial calculation and all endpoints}
\begin{proof}
For \(x,y>0\), differentiate the finite log expression on \((0,1)\).
Its first derivative vanishes exactly when
\(x(1-p)=yp\), hence \(p=x/(x+y)=x/n\).
The displayed second derivative is strictly negative. The derivative
is positive below that point and negative above it, so this is the
unique interior maximum. The likelihood is zero at both endpoints
and positive inside, making it the unique closed-interval maximum too.
If \(x=0,y=n\), the likelihood is \((1-p)^n\), strictly decreasing
on \([0,1]\); its unique maximum is at zero. If \(y=0,x=n\),
\(p^n\) is strictly increasing and its unique maximum is at one.
These exhaust the cases because \(n\geq1\). Removing the attaining
endpoint from the parameter space leaves a supremum but no maximum.
The empty-factor convention is the product of no factors, equal to one;
it does not change an impossible positive-count endpoint probability.

Each Bernoulli variable has mean \(p\) and variance \(p(1-p)\),
by its two masses. Independence makes the cross covariances zero,
so the sum has mean \(np\) and variance \(np(1-p)\).
Divide by \(n\) for the stated proportion moments. On the event
\(|\widehat p-p|>\epsilon\), its squared deviation exceeds
\(\epsilon^2\). Integrating that inequality bounds the event's
probability by its variance divided by \(\epsilon^2\). The bound
goes to zero for every fixed parameter and positive \(\epsilon\),
including the deterministic endpoints. This is convergence in probability.

An ordered string with \(x\) successes has probability
\(p^x(1-p)^{n-x}\). There are \(\binom nx\) such strings:
choose their success positions. Summing gives the count mass.
When that mass is positive, dividing the string mass by it gives
\(1/\binom nx\), independently of \(p\). Define this same uniform
kernel for every count, even for counts impossible at an endpoint;
the endpoint joint-law identity uses only its possible count.
Thus one parameter-free conditional kernel satisfies the definition
of sufficiency throughout the family. The map \(x\mapsto x/n\)
is injective at fixed positive \(n\), so conditioning on the proportion
is the same partition as conditioning on the count.
\end{proof}
\end{annexitem}

\begin{annexitem}{annex:p15.fe.coordinates}{thm:p15-fe-coordinates}{Coordinates, prior measures and density modes}
\begin{proof}
The map \(g(\theta)=(1+\sin\theta)/2\) is a continuous strictly
increasing bijection of the stated closed intervals. Thus its composition
with a likelihood attains its maximum exactly at the inverse image of
the likelihood's maximizing point; no parameter-density Jacobian appears
in this optimization. The preceding theorem gives
\(\sin\widehat\theta=2x/n-1=(x-y)/n\).
For interior \(x,y>0\), the transformed log derivative alternatively
gives
\(x/(1+\sin\theta)=y/(1-\sin\theta)\): multiply by the
positive denominators to obtain
\(x-y=(x+y)\sin\theta\). There is no extra factor two.
At endpoints use the bijection argument, not division by zero cosine.

In the open intervals,
\[
 \frac{dp}{d\theta}=\frac{\cos\theta}{2}
                    =\sqrt{p(1-p)},\qquad
 \frac{d\theta}{dp}=\frac1{\sqrt{p(1-p)}}.
\]
The ordinary change-of-variable identity therefore transports uniform
\(\theta\) density \(1/\pi\) to the stated arcsine density, and
uniform \(p\) density to \(\cos\theta/2\).
They integrate to one by the same substitution; endpoints have zero
mass and require no finite derivative there. They are different prior
measures, not two descriptions of one flat prior.

For one success and one failure the likelihood factor is \(p(1-p)\).
Its uniform-\(p\) normalizer is
\(\int_0^1(p-p^2)dp=1/2-1/3=1/6\), giving density \(6p(1-p)\).
For the arcsine prior, the normalizer is
\(\pi^{-1}\int_0^1\sqrt{p(1-p)}dp\).
Under the same sine substitution,
\[
 \int_0^1\sqrt{p(1-p)}dp
    =\frac14\int_{-\pi/2}^{\pi/2}\cos^2\theta\,d\theta
    =\frac\pi8.
\]
The last equality follows from
\(\cos^2\theta=(1+\cos2\theta)/2\).
The normalizer is consequently \(1/8\), giving the other stated density.
The two continuous densities differ, for example at \(p=1/2\),
where their values are \(3/2\) and \(4/\pi\).
They are both normalized by their own declared prior construction.

For \(x=1,y=2\), the uniform-\(p\) posterior is proportional
to \(p(1-p)^2\). Its positive interior log derivative is
\(1/p-2/(1-p)\), so its unique density mode is \(1/3\);
its density vanishes at both endpoints. Transport this very same
posterior measure to \(\theta\), multiplying its density by
\(dp/d\theta\). Expressed through the increasing coordinate \(p\),
the new height is proportional to \(p^{3/2}(1-p)^{5/2}\).
Its log derivative is \((3/2)/p-(5/2)/(1-p)\), yielding the
unique height maximum \(p=3/8\), again with zero endpoint heights.
The likelihood has no such Jacobian and remains maximized at \(1/3\).
This distinguishes change of density mode from change of measure and
from maximization of an observed-data likelihood.
\end{proof}
\end{annexitem}

\begin{annexitem}{annex:p15.fe.score}{thm:p15-fe-score}{Finite derivatives and the variance-information inequality}
\begin{proof}
Positivity makes each log derivative defined. Since the outcome set is
finite, differentiation of its normalization is ordinary differentiation
of a finite sum. Thus
\[
 E_\eta U_\eta=\sum_z q_\eta(z)\frac{q'_\eta(z)}{q_\eta(z)}
              =\sum_zq'_\eta(z)=0.
\]
The statistic is fixed independently of \(\eta\), so
\(m'(\eta)=\sum_zT(z)q'_\eta(z)=E_\eta(TU_\eta)\).
Subtracting \(m E_\eta U_\eta=0\) proves the covariance identity.
All moments in this fixed finite space are finite.
If \(J>0\), apply nonnegativity of the square to
\(T-m-cU_\eta\), with \(c=m'/J\). Expansion gives
\[
 0\leq E_\eta(T-m-cU_\eta)^2
        =\operatorname{Var}_\eta(T)-\frac{(m')^2}{J}.
\]
Multiplication by positive \(J\) proves the bound.
If \(J=0\), every positive-mass value of \(U_\eta\) is zero,
so the covariance identity instead gives \(m'=0\).

With twice differentiable masses,
\(\partial_\eta^2\log q=q''/q-(q'/q)^2\).
Average this expression: normalization differentiated twice makes
\(\sum_zq''_\eta(z)=0\), giving the curvature identity.
No integral, infinite sum or limiting operation was interchanged.

For \(n\) independent repetitions, the joint mass is the product
of the one-trial masses and its score is the sum of their scores.
Each has mean zero and second moment \(J\); independence makes every
cross expectation zero. The joint information is therefore \(nJ\).
The preceding covariance/square argument applies on the finite product
space to any fixed statistic of that product. If its expectation equals
\(\eta\) throughout the open interval, its mean derivative is one,
giving variance at least \(1/(nJ)\) when \(J>0\).
Pointwise unbiasedness at just one parameter would not justify that
derivative.

For Bernoulli masses, the scores at success/failure are
\(1/p\) and \(-1/(1-p)\), respectively. Hence
\[
 J=p/p^2+(1-p)/(1-p)^2=1/(p(1-p)).
\]
The exact proportion variance already proved is \(1/(nJ)\), so
attainment holds. Equivalently its centered value is
\((p(1-p)/n)\) times the joint score
\((K-np)/(p(1-p))\), making the square remainder zero.
These arguments assert nothing at zero-probability endpoint score values.
For the count likelihood, the negative second derivative is
\(x/p^2+y/(1-p)^2\). Its expectation under the true parameter is
\(np/p^2+n(1-p)/(1-p)^2=n/(p(1-p))\).
At an interior maximum substitute \(x=n\widehat p\) and
\(y=n(1-\widehat p)\) instead, obtaining
\(n/(\widehat p(1-\widehat p))\). This last substitution uses
the observed counts, not their expectation under the true parameter.
\end{proof}
\end{annexitem}

\begin{annexitem}{annex:p15.fe.failures}{thm:p15-fe-failures}{Three distinct finite countermodels}
\begin{proof}
On \([0,1]\), the three displayed masses are at least
\(1/5,3/10,1/5\), respectively, and their numerators sum to five.
They therefore define a positive common-support experiment.
Its first mass determines \(\eta^2=2-5q_\eta(1)\); since
\(\eta\geq0\), it determines \(\eta\). Thus the target is identified.
The first two masses strictly decrease as \(\eta\) increases from
zero; each has its unique maximum at zero. The third strictly increases
and has its unique maximum at one. This gives precisely the claimed
MLE statistic and its fibers.

Conditional on its fiber \(T=0\), the probability of outcome one is
\[
 \frac{q_\eta(1)}{q_\eta(1)+q_\eta(2)}
                =\frac{2-\eta^2}{4-3\eta^2/2}.
\]
The denominator is positive throughout. At zero it equals \(1/2\);
at one it equals \(1/(5/2)=2/5\). Thus a parameter-free conditional
kernel does not exist, and agreement of the likelihood maximizing point
inside that fiber did not imply sufficiency.

For the extension to \([-1,1]\), each mass depends only on
\(\eta^2\); the signed target is consequently unidentified.
For any \(0<a\leq1\), the ordered independent sample law at \(a\)
equals that at \(-a\), for every sample size. Suppose an estimator
sequence were consistent at both. The two events that its distance
from \(a\), respectively \(-a\), is less than \(a/2\) are
disjoint, so under their common law their probabilities sum to at most
one. Consistency would make both probabilities tend to one, a
contradiction. This argument concerns any statistic sequence, not just
the displayed one-observation MLE.

For the dependent binary experiment, define \(Z_i=B\) for every
\(i\). Its proportion is exactly \(B\), so its mean is \(p\)
and its variance \(p(1-p)\), without an \(n\) divisor.
Choose \(0<\epsilon<\min(p,1-p)\). Both possible values zero and
one have distance more than \(\epsilon\) from \(p\), giving
error probability one for every \(n\). Under independence, in
contrast, a mixed string has positive probability when \(n\geq2\)
and \(0<p<1\); under this experiment every mixed string has
probability zero. It is a different joint family, not the same
experiment with a better optimization algorithm. The variance and
consistency conclusions from the independent model cannot be transferred.
\end{proof}
\end{annexitem}

\begin{annexitem}{annex:p15.fe.cauchy}{thm:p15-fe-cauchy}{The Cauchy location law and finite median control}
\begin{proof}
Write \(X_i=m+Y_i\), where the \(Y_i\) are independent standard
Cauchy variables. The already supplied
\cref{thm:p15-characteristic-inverse-examples} and its full annex proof
give \(\sum_iY_i/n\) the standard Cauchy law. Thus
\(\overline X_n-m\) has that law. Integrating its density from
\(-\epsilon\) to \(\epsilon\), using the arctangent antiderivative,
gives \(2\arctan\epsilon/\pi\).
On \([1,R]\),
\[
 \int_1^R\frac{u}{\pi(1+u^2)}du
       =\frac1{2\pi}\log\frac{1+R^2}{2}\longrightarrow\infty.
\]
Both one-sided first moments have infinite absolute part. Also
\(u^2/(1+u^2)\geq1/2\) for \(u\geq1\), proving the infinite
second moment. Adding a fixed location does not restore absolute
integrability: \(|X|\geq|Y|-|m|\), and for large \(|Y|\),
\(X^2\geq Y^2/4\). Symmetric principal-value cancellation is
not the definition of a finite expectation.

The Cauchy distribution function is
\(F_m(t)=1/2+\arctan(t-m)/\pi\), directly from the same
antiderivative. Put \(\delta=\delta_\epsilon>0\).
The count of observations at most \(m+\epsilon\) is binomial
with parameter \(1/2+\delta\). If the middle order statistic
exceeds that threshold, this count is at most \(r\), while its
mean is \(n(1/2+\delta)\). The distance below its mean is then
at least \(n\delta\). Its variance is at most \(n/4\), so the
same square-deviation probability argument used for the binomial
proportion bounds this event by \(1/(4n\delta^2)\).

At \(m-\epsilon\), the count has mean \(n(1/2-\delta)\).
If the median is below this threshold, at least \(r+1\)
observations are at most the threshold, giving distance above the
mean at least \(n\delta\). The identical variance bound applies.
Continuity makes equality to either fixed threshold a probability-zero
event. Sum the two tail bounds and also use the trivial probability
bound one. For each fixed \(\epsilon>0\) this tends to zero,
proving median consistency along odd sample sizes. Neither a quantile
CLT nor an asymptotic MLE variance was assumed or proved here.
\end{proof}
\end{annexitem}
